\section{Empirical study design}\label{sec:design}

\subsection{Research questions}
\begin{description}[leftmargin=2.2em,style=nextline,itemsep=2pt]
\item[RQ1 (Reliability of detection).] How much does the detection quality of learning-based detectors drop when the evaluation moves from a random split to project-disjoint, temporal and cross-dataset conditions, and is AUROC a faithful predictor of operational value?
\item[RQ2 (Calibration).] How well-calibrated are the probabilities of these detectors, and how much do post-hoc recalibration methods help under shift?
\item[RQ3 (Guarantees under shift).] Does the conformal clearance rule keep its miss-rate budget when calibrated on source data, and which remedies (covariate-shift weighting, labelled target recalibration, the PAC variant) restore it? How many labelled target positives are needed?
\item[RQ4 (Accountability).] Is the guarantee distributed evenly across function sizes and vulnerability classes (CWEs), and does Mondrian calibration repair imbalances? How consistent are the triage decisions and the explanations under semantics-preserving rewrites and across datasets?
\item[RQ5 (Operational value).] How much reviewer effort does risk-controlled triage save at a valid miss budget, how does it compare with a rule-based static analyser, and what does it imply under an explicit cost model?
\item[RQ6 (Pre-trained code model).] Do the conclusions hold for a scorer built on a pre-trained code model (CodeBERT), and do such models exploit dataset artifacts that lexical models cannot see?
\end{description}

\subsection{Datasets and data-quality audit}
We use two public C/C++ function-level vulnerability datasets that differ in provenance.
\textbf{DiverseVul}~\cite{chen2023diversevul} contains functions collected from security-issue websites and the commits that fix them; each record carries the project, the fixing commit and (where available) the CWE.
\textbf{Big-Vul}~\cite{fan2020bigvul} links CVE records to fixing commits; we use the pre-fix function body as the vulnerable sample, the unchanged functions from the same commits as benign samples, and take the publication year from the CVE identifier as a temporal key.
Both were obtained from the Hugging Face hub (\texttt{claudios/DiverseVul}, \texttt{bstee615/bigvul}).
We deliberately did not use the Juliet test suite, whose synthetic flaws make random-split results saturate, nor PrimeVul, which was not retrievable in our environment.

\paragraph{Pre-processing and audit}
For each function we strip comments, collapse whitespace and hash the result. Within a dataset we remove exact duplicates, and we drop every hash that occurs with conflicting labels. Across datasets we compute the overlap of normalised hashes.
Table~\ref{tab:data} reports the audit; the consequences for evaluation are discussed in Section~\ref{sec:rq1}.

\begin{table}[t]
\centering\footnotesize
\caption{Dataset audit after normalisation (comment/whitespace stripping). 69,268 normalised functions occur in both datasets: 42.7\% of clean Big-Vul and 21.3\% of clean DiverseVul.}\label{tab:data}
\resizebox{\linewidth}{!}{\begin{tabular}{lrrrrrr}
\toprule
Dataset & Raw & Exact duplicates & Label-conflict & Clean & Vulnerable & Projects \\
\midrule
DiverseVul & 328,907 & 3,025 (0.9\%) & 1,497 & 325,138 & 17,689 (5.4\%) & 800 \\
Big-Vul & 216,568 & 53,373 (24.6\%) & 2,333 & 162,274 & 8,055 (5.0\%) & 309 \\
\bottomrule
\end{tabular}}
\end{table}


\subsection{Evaluation scenarios}
Table~\ref{tab:scen} defines seven scenarios. In every scenario the detector is trained on the \emph{train} part; the \emph{ID calibration} part is a random hold-out of the training distribution (what a practitioner would naturally use); the \emph{test} part is where all reported numbers are measured; and, for the shifted scenarios, a labelled \emph{target pool} drawn from the deployment distribution is available for the recalibration experiments but disjoint from the test part.
\begin{itemize}[itemsep=1pt]
\item \textbf{Random} (DV-R, BV-R): 70/15/15 split of deduplicated functions; calibration and test are exchangeable.
\item \textbf{Project} (DV-P, BV-P): 70\% of projects are used for training and calibration; the remaining projects form the test set (a new codebase is deployed). A random 25\% of the held-out projects' functions is the target pool.
\item \textbf{Temporal} (BV-T): train and ID calibration on CVEs published up to 2015, target pool = 2016, test = 2017--2019 (a model deployed on newer vulnerabilities).
\item \textbf{Cross-dataset} (DV$\to$BV, BV$\to$DV): train on 85\% of one dataset, calibrate on the remaining 15\%, test on the other dataset after \emph{removing every test function whose normalised hash occurs in the training data}; a random 25\% of the remaining target data is the pool. Removed (leaked) functions are evaluated separately to measure how much leakage would have inflated the result.
\end{itemize}
Each scenario is repeated with three seeds (different partitions and model initialisations).

\begin{table}[t]
\centering\footnotesize
\caption{Evaluation scenarios (seed 0 sizes; training keeps all vulnerable and at most 150\,000 benign functions).}\label{tab:scen}
\resizebox{\linewidth}{!}{\begin{tabular}{llrrrrr}
\toprule
Scenario & Shift & Train & ID cal. & Target pool & Test & Test pos. \\
\midrule
DV random & random 70/15/15 & 162,475 & 48,879 & -- & 48,553 & 2,637 \\
BV random & random 70/15/15 & 113,482 & 24,489 & -- & 24,303 & 1,254 \\
DV project & held-out projects (30\%) & 159,829 & 30,432 & 31,053 & 91,423 & 4,522 \\
BV project & held-out projects (30\%) & 124,551 & 22,227 & 3,826 & 11,670 & 699 \\
BV temporal & train $\le$2015; test 2017--19 & 60,101 & 10,850 & 27,022 & 52,669 & 2,288 \\
DV$\to$BV & train DV, test BV (leak-free) & 165,052 & 48,553 & 25,869 & 77,380 & 4,240 \\
BV$\to$DV & train BV, test DV (leak-free) & 137,971 & 24,303 & 66,735 & 199,348 & 11,207 \\
\bottomrule
\end{tabular}}
\end{table}


\subsection{Detectors}
Because the experiments must be fully reproducible on commodity CPUs and over 20 scenario--seed pairs, we use four lightweight, widely used detectors on the full corpora, complemented by a frozen-CodeBERT experiment on a subsample (Section~\ref{sec:design_cb}); the absence of fine-tuned code models is a limitation discussed in Section~\ref{sec:threats}, and the triage layer is detector-agnostic.
\begin{itemize}[itemsep=1pt]
\item \textbf{TF-IDF LR} and \textbf{TF-IDF SVM}: unigram+bigram code tokens hashed to $2^{19}$ features, sublinear TF-IDF, logistic regression ($C{=}20$) or linear SVM ($C{=}0.3$).
\item \textbf{LightGBM}~\cite{ke2017lightgbm}: gradient-boosted trees on 71 engineered code metrics (size, nesting depth, cyclomatic proxy, pointer/cast/bit-operation counts, counts of 32 security-relevant API calls) plus the 20\,000 most frequent hashed tokens; early stopping on 10\% of the training data.
\item \textbf{LR+LGBM}: sum of standardised LR margin and LightGBM logit.
\end{itemize}
Training sets keep all vulnerable functions and at most 150\,000 benign ones; this does not affect the triage layer, which depends on score ranks only.
The rule-based baseline is \textbf{Flawfinder} 2.0.20, with the highest risk level among a function's hits (then the number of hits) as score.

\subsection{Pre-trained code model experiment (RQ6)}\label{sec:design_cb}
To test whether the findings depend on the lightweight detectors, we add scorers built on CodeBERT~\cite{feng2020codebert} (\texttt{microsoft/codebert-base}). Because the study ran without a GPU, the encoder is \emph{frozen}: we embed the first 256 tokens of each function and use the concatenation of the [CLS] vector and the mean-pooled last layer (1\,536 dimensions) as features of a standardised logistic regression ($C{=}0.05$) and of a LightGBM (embedding plus the 71 code metrics); their standardised sum is the ``CodeBERT ensemble''. Embedding the full corpora (487\,000 functions) is infeasible on CPU, so we use a stratified subsample of each corpus with all scenarios, splits and seeds unchanged: 5\,000 vulnerable and 10\,000 benign functions per dataset (30\,000 in total; the benign functions are a uniform sample, so positive-class quantities such as the miss rate are unaffected, while AUPRC and SAR are weighted back to the true prevalence). All scorers in this experiment, including the TF-IDF models, are retrained on the subsample so that comparisons are matched.
Two input variants are evaluated: \emph{raw text}, and \emph{normalised text} (comments removed, whitespace collapsed), the same normalisation used in the main study. A \emph{formatting-only} LightGBM on 15 formatting statistics of the raw text (indentation, trailing whitespace, comment markers, line lengths) serves as a shortcut probe. For cross-dataset scenarios this experiment is stricter than the main study: any target function that occurs anywhere in the full source corpus is removed.

\subsection{Metrics}
\textbf{Detection:} AUROC, AUPRC, and F1, MCC, precision, recall at the F1-optimal threshold chosen on ID calibration data.
\textbf{Calibration:} expected calibration error with 15 equal-mass bins, $\mathrm{ECE}=\sum_b\frac{|B_b|}{n}\lvert\bar y_b-\bar p_b\rvert$, and the Brier score, for raw probabilities, Platt scaling and isotonic regression fitted on ID calibration data.
\textbf{Triage:} realised miss rate $M$ on test positives against the nominal $\alpha\in\{0.02,0.05,0.10,0.20\}$, the share of cleared functions (SAR, Eq.~\eqref{eq:sar}), the review effort ERE, and the violation frequency $\Pr(M>\alpha+0.02)$ across 200 random calibration draws of 300 positives (60 for the weighted variant); the $2$-point tolerance reflects the sampling error of the test-set miss rate.
For certified flagging we report the flagged share and the realised precision.
\textbf{Accountability:} miss rate per length tercile and per CWE.
\textbf{Robustness:} the share of triage decisions that change after identifier renaming, dead-code insertion, or both, and the miss rate on rewritten inputs.
\textbf{Explanation:} TreeSHAP~\cite{lundberg2020treeshap} deletion fidelity, rank stability and cross-dataset consistency of feature importances.
\textbf{Operational:} expected cost~\eqref{eq:cost} with the budget $\alpha^\star$ chosen on calibration data.

\subsection{Statistical protocol}
Results are reported as means over the three seeds with standard deviations. Detectors are compared with the Friedman test over the 21 scenario--seed blocks; paired comparisons with the best detector use the Wilcoxon signed-rank test and Cliff's $\delta$ as effect size.
Whether the realised miss rate exceeds its nominal level is tested with a one-sided Wilcoxon test on the per-block differences between realised and nominal miss rate. Relations between quantities (AUROC versus clearable share, KS distance versus excess miss rate) are measured with Spearman's $\rho$. All thresholds, hyper-parameters and splits are fixed before looking at test data; the 2-point tolerance and the budgets $\alpha$ are the only analysis choices and are reported in full.

\subsection{Implementation and availability}
All code (data preparation, feature extraction, models, triage layer, experiments, figure/table generation and this manuscript) is in the repository accompanying the paper, with a unit test that checks the clearance rule against its $p$-value definition and its finite-sample validity by simulation. Every number in the paper is produced by the scripts from cached detector scores.
